\section{Reproducibility and Setup}
\label{app:reproducibility}

\subsection{Hardware}

All results use fixed public LeWM checkpoints and no additional LeWM training. Local
development and most PushT analyses were run on the Apple Silicon workstation
recorded in the repository setup notes: macOS 26.3 ARM on an Apple M5 Pro, using
PyTorch MPS when available. The MPS device is visible from a normal host
terminal; the Codex sandbox can fail to expose Metal devices, so all MPS claims
refer to host-terminal runs. The RunPod closeout record reports an NVIDIA RTX
5090 server with CUDA 13.0 and 32GB VRAM for the final GPU-only checks. The
repository also includes \texttt{scripts/setup\_runpod.sh}, whose header targets
a fresh RunPod instance with system Python 3.12, CUDA 12.8, and pre-installed
\texttt{torch}/\texttt{torchvision}; the closeout log is the source for the
actual final RunPod runtime.

\begin{table}[h]
\centering
\scriptsize
\begin{tabular}{p{0.22\linewidth}p{0.28\linewidth}p{0.40\linewidth}}
\toprule
Hardware & Evidence in repo & Experiments run there \\
\midrule
Apple Silicon MPS workstation & \texttt{README.md}; \texttt{progress\_log.md}; \texttt{mps\_setup.md} & Phase 0 PushT/Cube baselines and sweeps; Track A 100-pair PushT artifact and three-cost pass; Stage 1A controls; Stage 1B PushT full projected CEM; PushT re-rank-only protocol matching; Cube Stage 1B re-rank-only; P2-0 learned-cost, DINOv2, and hybrid-CEM analyses. \\
RunPod RTX 5090 GPU server & \texttt{closeout\_session\_summary.md} & Cube anomaly multi-seed sanity; PushT Subspace-CEM Stage A; CUDA-side Cube full projected CEM/protocol matching artifacts. \\
CPU / offline aggregation & Artifact metadata and memos & Orthogonal sanity check, CEM-gap aggregation, taxonomy table, paper figure generation, and LaTeX compilation. \\
\bottomrule
\end{tabular}
\caption{Compute placement for the main audit artifacts.}
\label{tab:appendix-hardware}
\end{table}

\subsection{Software Environment}

The standard local setup command is \texttt{bash scripts/setup\_env.sh}. This
creates the \texttt{lewm-audit} conda environment with Python 3.11. The requested
search for environment YAML files returned no \texttt{*.yml} or \texttt{*.yaml}
environment file, so the repository should be treated as script-configured rather
than YAML-configured. The setup and progress logs verify
\texttt{stable-worldmodel==0.0.6}, editable
\texttt{stable-pretraining==0.1.7.dev44+gd8d8f568d.d20260501},
\texttt{ogbench==1.2.1}, and local PyTorch
\texttt{2.13.0.dev20260430}. RunPod final checks used miniconda,
\texttt{lewm-audit}, and PyTorch \texttt{2.11+cu128} according to the closeout
summary; the one-shot RunPod script installs \texttt{stable-worldmodel[train]},
\texttt{ogbench}, \texttt{gymnasium}, \texttt{opencv-python-headless},
\texttt{matplotlib}, \texttt{seaborn}, \texttt{scipy}, and
\texttt{scikit-learn} after confirming the pre-installed PyTorch package.

The only tracked local patch is \texttt{patches/lewm\_mps\_fix.patch}. It changes
the upstream LeWM \texttt{jepa.py} tensor-transfer path so that floating tensors
moved to an MPS device are cast to \texttt{torch.float32}; without this safeguard, MPS
rejects upstream \texttt{float64} tensors. Re-application instructions are in
\texttt{docs/mps\_setup.md}. The paper was compiled with \texttt{tectonic
main.tex} from \texttt{paper/}, using the ICLR 2025 template.

\begin{table}[h]
\centering
\scriptsize
\begin{tabular}{p{0.24\linewidth}p{0.32\linewidth}p{0.34\linewidth}}
\toprule
Component & Recorded version or setting & Source \\
\midrule
Conda environment & \texttt{lewm-audit} & \texttt{README.md}; setup scripts \\
Python & \texttt{3.11} local; system Python 3.12 targeted by RunPod script & \texttt{README.md}; setup scripts \\
\texttt{stable-worldmodel} & \texttt{0.0.6} local verified install & \texttt{README.md}; progress log \\
\texttt{stable-pretraining} & \begin{tabular}[t]{@{}l@{}}\texttt{0.1.7.dev44+}\\\texttt{gd8d8f568d.d20260501}\end{tabular} editable install & \texttt{README.md}; progress log \\
\texttt{ogbench} & \texttt{1.2.1} local verified install & \texttt{README.md}; progress log \\
PyTorch & \texttt{2.13.0.dev20260430} local MPS; \texttt{2.11+cu128} RunPod closeout & \texttt{README.md}; closeout summary \\
CUDA / GPU & CUDA 13.0, RTX 5090, 32GB VRAM for final RunPod runs & closeout summary \\
LaTeX & \texttt{tectonic main.tex} & closeout summary \\
\bottomrule
\end{tabular}
\caption{Software versions recorded in setup scripts and progress logs.}
\label{tab:appendix-software}
\end{table}

\subsection{Checkpoints}

Checkpoint download is handled by \texttt{scripts/download\_checkpoints.sh} on
local machines and by the equivalent checkpoint step in
\texttt{scripts/setup\_runpod.sh} on RunPod. Both scripts use the Hugging Face
CLI, with repository type \texttt{model}, to download the official public LeWM
checkpoint repositories. The root README links the LeWM Hugging Face collection
at \url{https://huggingface.co/collections/quentinll/lewm}. The exact Hugging
Face snapshot revision was not pinned at the time of the audit. All numerical
results are reproducible from the frozen JSON artifacts listed in
Table~\ref{tab:appendix-commits}.

\begin{table}[h]
\centering
\scriptsize
\setlength{\tabcolsep}{3pt}
\begin{tabular}{p{0.15\linewidth}p{0.21\linewidth}p{0.27\linewidth}p{0.25\linewidth}}
\toprule
Environment & Hugging Face model ID & Raw local files & Converted checkpoint \\
\midrule
PushT & \texttt{quentinll/lewm-pusht} & \begin{tabular}[t]{@{}l@{}}\texttt{weights.pt} and\\\texttt{config.json} under\\\texttt{checkpoints/lewm-pusht/}\end{tabular} & \begin{tabular}[t]{@{}l@{}}\texttt{lewm\_object.ckpt}\\under \texttt{checkpoints/}\\\texttt{converted/lewm-pusht/}\end{tabular} \\
OGBench-Cube & \texttt{quentinll/lewm-cube} & \begin{tabular}[t]{@{}l@{}}\texttt{weights.pt} and\\\texttt{config.json} under\\\texttt{checkpoints/lewm-cube/}\end{tabular} & \begin{tabular}[t]{@{}l@{}}\texttt{lewm\_object.ckpt}\\under \texttt{checkpoints/}\\\texttt{converted/lewm-cube/}\end{tabular} \\
\bottomrule
\end{tabular}
\caption{Official public LeWM checkpoints used by the audit.}
\label{tab:appendix-checkpoints}
\end{table}

\texttt{scripts/verify\_checkpoint.py} reconstructs the LeWM \texttt{JEPA} model
from \texttt{config.json}, loads \texttt{weights.pt}, saves a
\texttt{stable\_worldmodel.policy.AutoCostModel}-compatible object checkpoint as
\texttt{lewm\_object.ckpt}, and verifies load plus forward passes. The progress
log records that both PushT and Cube checkpoint files are about 69MB, that the
converted checkpoints produce encoder embeddings of shape \((1,192)\), and that
predictor verification produced shape \((1,3,192)\). The verified model has
about 18.0M parameters: 5.5M in the encoder and 12.5M in the predictor stack.
PushT uses raw action dimension 2 and effective checkpoint action dimension 10;
Cube uses raw action dimension 5 and effective checkpoint action dimension 25,
matching \texttt{action\_block=5}.

\subsection{Evaluation Seeds}

The audit uses fixed dataset-pair artifacts where possible and then varies only
the intended stochastic object. For Phase 0 baselines and offset sweeps, the
default command-line seed is 42. The later paper artifacts use seed 0 for fixed
pair selection, CEM sampling base seeds, and single-run controls unless a
multi-seed check is explicitly listed. Gaussian projection matrices have shape
\([192,m]\), entries \(\mathcal{N}(0,1)/\sqrt{m}\), and are fixed per
\((m,\texttt{projection\_seed})\) across pairs and CEM iterations.

\begin{table}[h]
\centering
\small
\begin{tabular}{p{0.28\linewidth}p{0.28\linewidth}p{0.34\linewidth}}
\toprule
Artifact or phase & Seeds & Notes \\
\midrule
Phase 0 PushT/Cube baselines and offset sweeps & 42 & Defaults in \texttt{eval\_pusht\_baseline.py}, \texttt{eval\_pusht\_sweep.py}, \texttt{eval\_cube\_baseline.py}, and \texttt{eval\_cube\_sweep.py}. \\
Track A PushT pair set and three-cost artifact & 0 & \texttt{track\_a\_pairs.json} and \texttt{track\_a\_three\_cost.json}; 100 pairs, 80 actions per pair. \\
Stage 1A C1-C5 controls & 0--9 & Ten-seed orthogonal, Gaussian projection, coordinate subset, Gaussian null, and shuffle controls. \\
Stage 1A C0/C6/C7 & 0 or deterministic single run & C0 is trained LeWM; C6 is randomly initialized LeWM seed 0; C7 DINOv2 feature extraction records seed 0. \\
Stage 1B PushT full projected CEM & base seed 0; projection seeds 0,1,2 & 63-pair smoke/full record; later full artifact uses dimensions 1--192 and the same projection seed set. \\
PushT re-rank-only protocol matching & base seed 0; projection seeds 0,1,2 & 100 pairs, 9 projection dimensions, 2700 projected records. \\
Cube Stage 1B re-rank-only & base seed 0; projection seeds 0,1,2 & 100 Cube pairs, 9 projection dimensions, 2700 projected records. \\
Cube full projected CEM smoke & base seed 0; projection seed 0 & 25 selected Cube pairs, dimensions \(\{1,8,32,64,192\}\). \\
Cube anomaly sanity & projection seeds 0,1,2 & Seed 0 from \texttt{cube\_full\_proj\_cem.json}; added seeds 1 and 2 in \texttt{cube\_anomaly\_sanity.json}. \\
P2-0 learned costs and hybrid CEM & 0 & Cost-head splits, planning evaluation, and oracle-budget sweeps record seed 0. \\
PushT Subspace-CEM Stage A & 0,1 & 30 pairs, \texttt{m=64} search with full-dimensional re-rank. \\
\bottomrule
\end{tabular}
\caption{Seed policy by experimental phase.}
\label{tab:appendix-seeds}
\end{table}

\subsection{Statistical Uncertainty}

Table~\ref{tab:appendix-uncertainty} collects the uncertainty intervals added
for the second-round revision.  The bootstrap intervals use the pair or
paired-pair experimental unit rather than individual candidates, and binary
repair outcomes use exact binomial intervals.

% Generated by scripts/revision/appendix_uncertainty_table.py
\begin{table*}[p]
\centering
\scriptsize
\setlength{\tabcolsep}{2.5pt}
\resizebox{\textwidth}{!}{%
\begin{tabular}{@{}L{0.18\textwidth}L{0.07\textwidth}L{0.20\textwidth}L{0.05\textwidth}L{0.06\textwidth}L{0.15\textwidth}L{0.12\textwidth}L{0.15\textwidth}@{}}
\toprule
Metric & Env. & Protocol / experiment & $N$ & Seeds & Point estimate & 95\% CI & Method \\
\midrule
$R_\mathrm{endpoint}$ ($m=64$) & PushT & Stage 1A C2 endpoint projection & 100 & 10 proj. & 0.473 & [0.395, 0.543] & pair bootstrap \\
$R_\mathrm{endpoint}$ ($m=192$) & PushT & Stage 1A C2 endpoint projection & 100 & 10 proj. & 0.495 & [0.411, 0.570] & pair bootstrap \\
$R_\mathrm{pool}(C_\mathrm{model})$ & PushT & Default CEM final pools & 100 & 1 & 0.084 & [0.037, 0.131] & pair bootstrap \\
$R_\mathrm{pool}(C_\mathrm{V3})$ & PushT & Actual-terminal encoder ranker & 100 & 1 & 0.258 & [0.146, 0.365] & pair bootstrap \\
$R_\mathrm{pool}(C_\mathrm{V1})$ & PushT & Physical hinge ranker & 100 & 1 & 0.666 & [0.595, 0.735] & pair bootstrap \\
$\Delta_\mathrm{CEM}$ & PushT & C0 endpoint ref. minus default pool & 100 & 1 & 0.422 & [0.374, 0.468] & pair bootstrap \\
$R_\mathrm{pool}(C_\mathrm{model})$ & Cube & Full projected CEM, $m=64$ & 50 & 3 proj. & -0.001 & [-0.018, 0.016] & pair bootstrap \\
$\Delta_\mathrm{CEM}$ & Cube & Endpoint $m=64$ minus full CEM pool & 50 & 3 proj. & 0.576 & [0.559, 0.594] & pair bootstrap \\
\addlinespace
$R_\mathrm{pool}(C_\mathrm{model})$ diff. & PushT & MPPI minus CEM & 30 & 3 & 0.079 (0.115 to 0.194) & [-0.044, 0.204] & paired bootstrap \\
Planning success diff. & PushT & MPPI minus CEM & 30 & 3 & -14.4 pp (53.3\% to 38.9\%) & [-26.7, -2.2] pp & paired bootstrap \\
Pool $C_\mathrm{real}$ std diff. & PushT & MPPI minus CEM & 30 & 3 & 21.530 (7.947 to 29.478) & [17.151, 25.854] & paired bootstrap \\
MPPI attribution fraction & PushT & CEM-to-MPPI gap fraction & 30 & 3 & 22.225\% & [-15.736, 48.094]\% & paired bootstrap \\
\addlinespace
Cost head Split 3 success & PushT & Learned cost repair & 16 & 1 & 0/16 (0.000\%) & [0.000, 20.591]\% & Clopper-Pearson \\
Cost head CEM-aware success & PushT & Learned cost repair & 16 & 1 & 1/16 (6.250\%) & [0.158, 30.232]\% & Clopper-Pearson \\
DINOv2 gap vs. LeWM & PushT & Track B endpoint replacement & 100 & 1 & 0.242 (0.261 vs. 0.503) & [0.115, 0.368] & pair bootstrap \\
Subspace-CEM success delta & PushT & Subspace minus default CEM & 30 & 2 & -8.3 pp (56.7\% to 48.3\%) & [-18.3, 0.0] pp & paired bootstrap \\
Hybrid CEM 20\% threshold & PushT & V1 oracle re-rank, $K=60$ & 16 & 1 & 7/16 (43.750\%) & [19.753, 70.122]\% & Clopper-Pearson \\
\bottomrule
\end{tabular}%
}
\caption{Uncertainty summary for the main endpoint-pool, optimizer, and repair metrics. Bootstrap intervals use 10{,}000 resamples over the experimental unit shown in the Method column. Exact intervals are Clopper--Pearson binomial intervals. Endpoint global Spearman intervals bootstrap pair clusters after rank-transforming endpoint records, so the point estimates match the global Spearman values reported in the main tables.}
\label{tab:appendix-uncertainty}
\end{table*}


\subsection{Reproducibility Notes}

The main fixed artifacts record the commit that generated them. The repository
history has continued after these pre-registration and closeout points, so the
artifact-level commit hashes are the reproducibility anchors rather than the
current branch head alone.

\begin{table}[h]
\centering
\scriptsize
\begin{tabular}{p{0.30\linewidth}p{0.22\linewidth}p{0.38\linewidth}}
\toprule
Anchor & Commit & Evidence \\
\midrule
Track A stratified pair artifact & \texttt{3ac6afa93c98} & \texttt{track\_a\_pairs.json} \\
Track A three-cost pass & \texttt{45d65afc1546} & \texttt{track\_a\_three\_cost.json} \\
Stage 1B smoke / full projected CEM & \texttt{7848d5b00f28}, \texttt{a48f9f3d0e18} & \texttt{stage1b\_smoke\_memo.md}; \texttt{stage1b\_full.json} \\
PushT re-rank-only and Cube Stage 1B & \texttt{a3388de9d6dc} & \texttt{pusht\_rerank\_only.json}; \texttt{cube\_stage1b.json} \\
Cube full projected CEM & \texttt{485b3390154f} & \texttt{cube\_full\_proj\_cem.json} \\
Protocol-matching memo & \texttt{7917feb3996e} & \texttt{protocol\_matching\_memo.md} \\
PushT Subspace-CEM Stage A & \texttt{adc573ac8c2b} & \texttt{stage\_a\_pusht.json} \\
\bottomrule
\end{tabular}
\caption{Commit anchors recorded by the reproducibility artifacts.}
\label{tab:appendix-commits}
\end{table}

Known non-determinism comes from stochastic CEM sampling, simulator rollouts,
random projections, GPU kernels, and backend differences between MPS and CUDA.
The code sets Python, NumPy, and PyTorch seeds in the Phase 2 evaluators where
the scripts expose a seed; CUDA scripts call \texttt{torch.cuda.manual\_seed\_all}
when CUDA is available. Some finite-sample baseline variation is expected: for
example, the reproduced PushT baseline is 98\% on 50 episodes versus the paper
reference point of 96\%, and Cube baseline is 66\% on 50 episodes versus the
paper reference point of 74\%.

Runtime estimates recorded in the repository are as follows. PushT MPS baseline
evaluation takes about 295 seconds for 50 episodes, with a one-time first-solve
shader warmup. The PushT offset sweep took about 1 hour on MPS. The Track A
100-pair three-cost run took about 9 minutes from its recorded start and finish
timestamps. PushT re-rank-only and Cube full projected CEM each report about
5 minutes 16 seconds. Cube Stage 1B re-rank-only reports about 1 hour 19
minutes. Final RunPod closeout runs report 7 minutes 8 seconds for the Cube
anomaly sanity check and 4 minutes 23 seconds for PushT Subspace-CEM Stage A.
